% Fichas de pesquisa: ATV, GER
% Conferido na web em 23/09/2026. "(a confirmar)" marca o que não foi verificado em fonte.

\section{Ativos e geração renovável}

\subsection*{Bases comuns}
\begin{itemize}
  \item \textbf{ONS Dados Abertos} (CC-BY, também espelhado no AWS) [8]: geração
    por usina em base horária [5], fator de capacidade horário de eólicas e
    solares [6], constrained-off semi-horário por usina desde 2021 [1--4],
    disponibilidade de linhas e transformadores [7]. O ONS avisa que os dados
    são revisados depois de publicados.
  \item \textbf{ANEEL}: SIGA (cadastro de usinas por CEG) [9], SIGEL (usinas e
    linhas em shapefile, aerogeradores no SIGEL-EOL) [10], BDGD [11],
    interrupções nas redes de distribuição [12].
  \item \textbf{Chave de junção}: o constrained-off traz \texttt{ceg} e
    \texttt{id\_ons} [4], o que liga ONS, SIGA e SIGEL pelo CEG; conjuntos de
    usinas aparecem sem CEG e exigem uma tabela de correspondência própria.
\end{itemize}

% ---------------------------------------------------------------------------
\subsection{ATV \textperiodcentered{} Ativos e manutenção}

\begin{ficha}{ATV-01 \textperiodcentered{} Índice de saúde de transformadores de potência}{V3}
\campo{Dados} Do cliente: cromatografia de gases dissolvidos (H$_2$, CH$_4$,
C$_2$H$_6$, C$_2$H$_4$, C$_2$H$_2$, CO, CO$_2$), físico-química do óleo,
furanos, carregamento, idade, histórico de ordens de serviço. Normas: IEC
60599:2022 [28]; IEEE C57.104-2019, que trocou limiares de gás total por
percentis por gás, deltas entre amostras e taxas anuais [27]; triângulos e
pentágonos de Duval [30]; índices de avaliação do CIGRE TB 761 [31]. Dados
públicos de referência: IEC TC10, com 117 amostras rotuladas [29].
\campo{Abordagens} Linha de base: regras normativas combinadas num índice
ponderado no estilo CIGRE. Recomendada: gradient boosting que usa as normas
como variáveis (razões, coordenadas de Duval, deltas, taxas, perda de vida
pelo IEEE C57.91 [33]), tendo como rótulo o evento posterior (falha ou
intervenção) e não a classe de Duval; com calibração e SHAP. Revisão recente:
Zahra et al.\ (2025) [32].
\campo{Métricas} F1 por classe de falha (diagnóstico); AUC-PR para falha ou
intervenção em 12 meses e acerto no topo da lista (prognóstico); concordância
com especialista.
\campo{Armadilhas} Validar só no IEC TC10 (117 amostras) dá confiança falsa;
taxas exigem três ou mais amostras, raras na prática; troca de laboratório e
tratamento do óleo criam saltos artificiais.
\end{ficha}

\begin{ficha}{ATV-02 \textperiodcentered{} Risco de falha por equipamento}{V3}
\campo{Dados} Públicos: BDGD (segmentos, transformadores, chaves, data de
instalação) [11] e interrupções por evento [12]; ONS para a rede básica [7].
Do cliente: cadastro com trocas, ordens de serviço estruturadas (ATV-05).
\campo{Abordagens} Linha de base: Weibull por classe e Cox. Recomendada:
random survival forest ou boosting com perda de sobrevivência e covariáveis
que mudam no tempo (clima, carregamento), tratando censura e truncamento à
esquerda [34,35]. Estado da arte: DeepSurv e modelos neurais de tempo
discreto [36]. No nível do alimentador: contagem de falhas por km e ano.
\campo{Métricas} C-index (inclusive dependente do tempo), Integrated Brier
Score, calibração em 12 meses, acerto no topo da lista.
\campo{Armadilhas} Truncamento à esquerda (viés de sobrevivente); trocas sem
registro; causas externas misturadas com desgaste (riscos competitivos).
\end{ficha}

\begin{ficha}{ATV-03 \textperiodcentered{} Vegetação próxima à rede}{V3}
\campo{Dados} Geometria da rede (BDGD [11], SIGEL [10]). Altura de dossel
pública: mapa global de 10~m para 2020 (Lang et al.\ 2023) [37] e de 1~m
(Meta e WRI, Tolan et al.\ 2024, imagens de 2017 a 2020) [38]. Sentinel-2
para atualização. LiDAR do cliente quando houver.
\campo{Abordagens} Linha de base: altura de dossel e NDVI cruzados com um
buffer do vão. Recomendada: segmentação de vegetação densa e regressão de
altura com Sentinel-2 e GEDI; caso brasileiro (SENAI e Cemig, 2025/2026) com
RMSE de 7,7~m e MAE de 5~m contra LiDAR [39]. Com LiDAR: distância direta
entre condutor e vegetação. Estado da arte: SAM adaptado para corredores de
transmissão [40].
\campo{Métricas} IoU e F1 da segmentação; erro de altura contra LiDAR; recall
de vãos críticos em campo; correlação com interrupções por vegetação [12].
\campo{Armadilhas} Erro de 5 a 8~m é da ordem da faixa de segurança: serve
para priorizar, não para medir distância. Mapas estáticos envelhecem com
crescimento e poda.
\end{ficha}

\begin{ficha}{ATV-04 \textperiodcentered{} Defeitos em imagens de inspeção}{V4}
\campo{Dados} Imagens de drone do cliente. Públicos: CPLID (defeitos
sintéticos) [21]; InsPLAD, com 10.607 imagens, 17 classes de ativo e 6 tipos
de defeito, mas só 402 amostras de defeito [22]; TTPLA (torres e cabos) [23];
UPID [24].
\campo{Abordagens} Linha de base: YOLO ajustado para componentes [26].
Recomendada: dois estágios, detectar o componente e classificar o defeito ou
detectar anomalia no recorte [22]. Estado da arte: RT-DETR (CVPR 2024) [25];
SAM para pré-anotação com humano no ciclo.
\campo{Métricas} mAP por classe; recall por defeito crítico; precisão por
estrutura (o que vira ordem de serviço).
\campo{Armadilhas} Defeitos raros; sintéticos não generalizam; tipos de
isolador brasileiros diferem; cada imagem precisa ser associada à estrutura
da BDGD.
\end{ficha}

\begin{ficha}{ATV-05 \textperiodcentered{} Estruturação de ordens de serviço e causas}{V3}
\campo{Dados} Texto livre das ordens de serviço do cliente; causas do dataset
de interrupções como rótulo fraco [12]. Esquemas: MaintIE (224 entidades, 6
relações) [41] e ISO 14224 para modos de falha [42].
\campo{Abordagens} Na classificação de modo de falha em 22 códigos ISO 14224,
um GPT-3.5 ajustado chegou a F1 0,80, contra 0,60 de um classificador de
texto e 0,46 do modelo sem ajuste [42]. Recomendada: modelo de linguagem com
saída JSON restrita a um esquema (componente, modo de falha, causa, ação),
exemplos do cliente, ajuste leve e revisão humana por amostragem.
\campo{Métricas} F1 por campo; taxa de JSON válido; ganho de C-index em ATV-02
com as variáveis extraídas.
\campo{Armadilhas} Os benchmarks são em inglês e o texto das ordens é curto,
abreviado e regional; cada campo extraído precisa guardar o trecho que o
justifica; LGPD.
\end{ficha}

\begin{ficha}{ATV-06 \textperiodcentered{} Vida útil remanescente e plano de manutenção}{V4}
\campo{Dados} ATV-01 e ATV-02; custos, equipes e orçamento do cliente;
consequência da falha pelos clientes afetados (BDGD, interrupções).
\campo{Abordagens} Linha de base: matriz risco × criticidade. Recomendada:
risco monetizado por ativo e otimização inteira sob equipes e regiões (a
confirmar como padrão). Estado da arte: otimização estocástica sob a
incerteza da curva de sobrevivência.
\campo{Armadilhas} Poucos fins de vida para validar; a otimização amplifica
erros de calibração.
\end{ficha}

% ---------------------------------------------------------------------------
\subsection{GER \textperiodcentered{} Geração renovável}

\begin{ficha}{GER-01 \textperiodcentered{} Geração prevista por parque (P10--P90)}{V2}
\campo{Dados} Localização e ficha: SIGA e SIGEL [9,10]. \textbf{Geração
observada pública}: ONS por usina, horária [5], e geração verificada
semi-horária no constrained-off [4]; cobre usinas Tipo I, II-B e II-C (II-C
agregadas em conjuntos). Previsão do tempo: ECMWF abriu todo o catálogo em
tempo real sob CC-BY em outubro de 2025 [53]. SCADA do cliente em médias de
10 minutos.
\campo{Abordagens} Linha de base: persistência, climatologia, curva de
potência aplicada à previsão. Recomendada: gradient boosting quantílico sobre
variáveis da previsão numérica, o método vencedor do GEFCom2014 [52], com
recalibração conformal. Onboarding sem histórico: foundation models com a
previsão do tempo como covariável futura (Chronos-2 [46], TimesFM 2.5 [47],
Moirai 2.0 [48]); histórias sintéticas geradas por modelo físico reduziram o
erro de 1,7 a 2 vezes em 440 usinas fotovoltaicas [50]. Nenhuma abordagem
domina sempre [49]. No Brasil, ERA5 superou o MERRA-2 para vento [54].
\campo{Métricas} nMAE e nRMSE pela capacidade; pinball e CRPS; cobertura
P10--P90; skill contra persistência; por horizonte e sem curtailment.
\campo{Armadilhas} \textbf{Treinar sobre a geração de referência, não sobre a
entregue}, para não aprender o curtailment [4]; revisões do ONS; vazamento
por tempo observado; UTC contra horário local.
\end{ficha}

\begin{ficha}{GER-02 \textperiodcentered{} Alertas de rampa}{V2}
\campo{Abordagens} Não há definição padrão: rampa é magnitude, duração
(tipicamente 1 a 4~h), instante e direção, com limiar pelo porte [55].
Linha de base: limiar de variação sobre a P50. Recomendada: probabilidade de
rampa calculada sobre \textbf{cenários coerentes no tempo}, não sobre quantis
marginais, que criam rampas espúrias.
\campo{Métricas} POD, FAR e CSI com tolerância de fase de uma hora; Brier;
antecedência; alertas por semana.
\end{ficha}

\begin{ficha}{GER-03 \textperiodcentered{} Nowcasting por satélite (0 a 6 h)}{V3}
\campo{Dados} GOES-19, GOES-East operacional desde 04/04/2025, no AWS sem
conta [13,14]; ABI com visível a 0,5~km e full disk a cada 10 minutos [15];
irradiância na superfície (DSR) a 2~km e 10 minutos no algoritmo Enterprise
[16]; radiação diária do INPE [17].
\campo{Abordagens} Linha de base: persistência do índice de céu claro e
vetores de movimento de nuvem. Recomendada: U-Net sobre sequências de índice
de nuvem, convertida em potência por usina [57,60]; caso brasileiro com
GOES-16 em Petrolina [59]. Estado da arte probabilístico: SHADECast, que
estende o horizonte útil em cerca de 26 minutos [56]. Acima de 2 a 3~h,
combinar com previsão numérica [53].
\campo{Métricas} RMSE de irradiância e potência; skill contra persistência
por antecedência; CRPS.
\campo{Armadilhas} Paralaxe e resolução contra usinas pequenas; latência;
descontinuidade na troca do GOES-16 para o GOES-19.
\end{ficha}

\begin{ficha}{GER-04 \textperiodcentered{} Curtailment detectado e marcado}{V2}
\campo{Dados} Constrained-off do ONS, dicionário v1.5 (09/2026): geração,
geração limitada, disponibilidade, referência, referência final (só para
REL, enviada à CCEE), razão (REL, CNF, ENE, PAR), origem (LOC, SIS), energia
não realizada e minutos por razão [4]. Relatório do GT Curtailment do ONS
[62]. SCADA com setpoints.
\campo{Abordagens} Linha de base: regras sobre a geração limitada e os
minutos de restrição; no SCADA, setpoint abaixo da nominal. Recomendada:
modelo de \textbf{potência disponível} (comportamento normal, GER-05);
curtailment é disponível menos entregue quando há limite, o que detecta
cortes não informados e contesta a referência do ONS; limpeza da curva de
potência por detecção de anomalias [64].
\campo{Armadilhas} Resolução de 30 minutos esconde cortes curtos; regras de
ressarcimento mudam por resolução (a confirmar).
\end{ficha}

\begin{ficha}{GER-05 \textperiodcentered{} Desempenho real × esperado}{V3}
\campo{Dados} SCADA, recurso, curtailment de GER-04. Públicos para
desenvolver: Kelmarsh e Penmanshiel (SCADA de 10 minutos, CC-BY) [18,19];
PVDAQ do NREL [68].
\campo{Abordagens} Linha de base: curva de potência por intervalos e
performance ratio. Recomendada: modelo de comportamento normal (boosting ou
GAM) treinado em período saudável; para fotovoltaica, pvlib ajustado.
Ferramentas: OpenOA (NREL) [65] e RdTools para degradação [71].
\campo{Armadilhas} Treinar o esperado sobre dados que já contêm a perda; a
limpeza pode reter só cerca de 31\% dos dados [64].
\end{ficha}

\begin{ficha}{GER-06 \textperiodcentered{} Saúde de turbinas}{V3}
\campo{Dados} SCADA com temperaturas de rolamentos, gerador e caixa
multiplicadora. Rotulados: CARE to Compare (Fraunhofer, 2024), 95 datasets, 36
turbinas, 44 eventos rotulados, com a métrica CARE (cobertura, acurácia,
confiabilidade, antecedência) [20].
\campo{Abordagens} Revisão: Tautz-Weinert e Watson (2017) [66]. Chesterman et
al.\ (2023), em 5 parques: Elastic Net competitivo com LightGBM e redes;
normalizar pela mediana da frota; excluir 4 meses antes da falha do treino;
nenhuma família domina [67]. Recomendada: comportamento normal com
normalização pela frota e alarme com persistência.
\campo{Métricas} CARE score; dias de antecedência; falsos alarmes por turbina
por ano.
\campo{Armadilhas} Poucas falhas; dados pré-falha contaminam o normal; clima
do Nordeste diferente do dos datasets europeus.
\end{ficha}

\begin{ficha}{GER-07 \textperiodcentered{} Sujidade e programação de limpeza}{V4}
\campo{Abordagens} Kimber (acúmulo linear até chuva acima de limiar) e HSU
(por massa de particulado), ambos no pvlib [69,70]; a partir da produção, SRR
(Deceglie et al.\ 2018) no RdTools [71]; relatório IEA PVPS T13-21:2022 [72].
Limiares de chuva que limpa variam de 0,3 a 20~mm/dia [73]. Decisão: limpar
quando a perda esperada até a próxima chuva prevista supera o custo.
\campo{Armadilhas} Confundir sujidade com degradação ou sombreamento;
limpezas não registradas quebram o SRR.
\end{ficha}

\begin{ficha}{GER-08 \textperiodcentered{} Geração agregada por submercado}{V3}
\campo{Abordagens} Soma das previsões por usina com upscaling para as sem
previsão; reconciliação hierárquica (usina, estado, submercado) [74] e versão
orientada a valor [75]. Quantis agregados exigem cenários conjuntos: soma de
quantis não é o quantil da soma.
\campo{Armadilhas} Entrada de novas usinas (atualizar pelo SIGA); dizer se a
previsão é de potência disponível ou entregue; GD invisível.
\end{ficha}

\subsection*{Referências desta seção}
{\small
\begin{enumerate}[label={[\arabic*]},leftmargin=2.2em,itemsep=1pt]
\item ONS, constrained-off de eólicas: \url{https://dados.ons.org.br/dataset/restricao_coff_eolica_usi}
\item ONS, constrained-off de eólicas por usina: \url{https://dados.ons.org.br/dataset/restricao_coff_eolica_detail}
\item ONS, constrained-off de fotovoltaicas: \url{https://dados.ons.org.br/dataset/restricao_coff_fotovoltaica_detail}
\item ONS, dicionário do constrained-off v1.5 (09/2026)
\item ONS, geração por usina horária: \url{https://dados.ons.org.br/dataset/geracao-usina-2}
\item ONS, fator de capacidade: \url{https://dados.ons.org.br/dataset/fator-capacidade-2}
\item ONS, disponibilidade da função transmissão: \url{https://dados.ons.org.br/dataset/ind_disponibilidade_ft_trlt}
\item ONS no AWS: \url{https://registry.opendata.aws/ons-opendata-portal/}
\item ANEEL, SIGA: \url{https://dadosabertos.aneel.gov.br/dataset/siga-sistema-de-informacoes-de-geracao-da-aneel}
\item ANEEL, SIGEL: \url{https://sigel.aneel.gov.br/portal/home/}
\item ANEEL, BDGD: \url{https://dadosabertos.aneel.gov.br/dataset/base-de-dados-geografica-da-distribuidora-bdgd}
\item ANEEL, interrupções: \url{https://dadosabertos.aneel.gov.br/dataset/interrupcoes-de-energia-eletrica-nas-redes-de-distribuicao}
\item NOAA GOES no AWS: \url{https://registry.opendata.aws/noaa-goes/}
\item NOAA, GOES-19 no NODD
\item GOES-R ABI: \url{https://www.goes-r.gov/spacesegment/abi.html}
\item NCEI, ABI L2 DSR
\item INPE, radiação solar média diária (GOES-16)
\item Kelmarsh: \url{https://zenodo.org/records/5841834}
\item Penmanshiel: \url{https://zenodo.org/records/5946808}
\item Gück, Roelofs e Faulstich (2024), CARE to Compare: \url{https://arxiv.org/abs/2404.10320}
\item CPLID: \url{https://github.com/InsulatorData/InsulatorDataSet}
\item Vieira e Silva et al.\ (2023), InsPLAD, \emph{IJRS} 44(23): \url{https://arxiv.org/abs/2311.01619}
\item Abdelfattah et al.\ (2020), TTPLA: \url{https://arxiv.org/abs/2010.10032}
\item UPID: \url{https://github.com/heitorcfelix/public-insulator-datasets}
\item Zhao et al.\ (2024), RT-DETR, CVPR: \url{https://arxiv.org/abs/2304.08069}
\item MAP-YOLOv8, \emph{Scientific Reports} (2025)
\item IEEE C57.104-2019
\item IEC 60599:2022
\item IEC TC10 em \emph{IET GTD} (2021); IEEE DataPort DGA dataset
\item Duval, triângulos e pentágonos
\item CIGRE TB 761 (2019)
\item Zahra et al.\ (2025): \url{https://arxiv.org/abs/2504.15310}
\item IEEE C57.91, guia de carregamento
\item Modelos de sobrevivência em redes de energia, \emph{Sensors} (2026), doi:10.3390/s26061915
\item Ishwaran et al.\ (2008), Random Survival Forests, \emph{Ann.\ Appl.\ Stat.}
\item Katzman et al., DeepSurv: \url{https://arxiv.org/abs/1606.00931}
\item Lang et al.\ (2023), \emph{Nature Ecology \& Evolution} 7:1778--1789
\item Tolan et al.\ (2024), \emph{RSE} 300:113888: \url{https://registry.opendata.aws/dataforgood-fb-forests/}
\item HVAS (SENAI e Cemig), ISPRS Annals X-3-W4-2025
\item SAM para corredores de transmissão: \url{https://arxiv.org/abs/2505.22105}
\item Bikaun et al.\ (2024), MaintIE, LREC-COLING
\item Stewart et al., LLMs para modos de falha: \url{https://arxiv.org/abs/2309.08181}
\item Benchmark de LLMs em logs de turbinas: \url{https://arxiv.org/abs/2509.06813}
\item Extração de KPIs de ordens de serviço: \url{https://arxiv.org/abs/2311.04064}
\item Brundage et al., Technical Language Processing, \emph{Manufacturing Letters} 27
\item Ansari et al., Chronos-2: \url{https://arxiv.org/abs/2510.15821}
\item TimesFM 2.5: \url{https://huggingface.co/google/timesfm-2.5-200m-pytorch}
\item Moirai 2.0: \url{https://arxiv.org/abs/2511.11698}
\item Za'ter e Hodge (2026): \url{https://arxiv.org/abs/2604.22077}
\item Longarini et al.\ (2026), histórias sintéticas: \url{https://arxiv.org/abs/2606.07457}
\item Chronos, TimesFM e Moirai, versões originais
\item Hong et al.\ (2016), GEFCom2014, \emph{IJF}; Landry et al., GBM
\item ECMWF, catálogo em tempo real aberto (2025)
\item Gruber et al., reanálises no Brasil: \url{https://arxiv.org/abs/1904.13083}
\item Gallego-Castillo et al.\ (2015), rampas, \emph{RSER} 52
\item Carpentieri et al., SHADECast: \url{https://arxiv.org/abs/2312.11966}
\item Straub et al.\ (2024), \emph{Solar RRL}
\item SunCast: \url{https://arxiv.org/abs/2201.06173}
\item GHI e DNI com GOES-16 em Petrolina, \emph{Int.\ J.\ Energy Environ.\ Eng.}
\item Nowcasting por satélite com deep learning, \emph{Solar Energy} (2024)
\item Shi et al., ConvLSTM: \url{https://arxiv.org/abs/1506.04214}
\item ONS, RT DGL-ONS 0189-2025, GT Curtailment
\item Canal Solar, cortes em 2026 (fonte secundária)
\item Limpeza de curva de potência em SCADA, \emph{Renewable Energy} (2022)
\item NREL OpenOA: \url{https://github.com/NREL/OpenOA}
\item Tautz-Weinert e Watson (2017), \emph{IET RPG} 11, doi:10.1049/iet-rpg.2016.0248
\item Chesterman et al.\ (2023), \emph{Wind Energy Science} 8:893
\item NREL PVDAQ: \url{https://data.openei.org/submissions/4568}
\item Modelo de sujidade de Kimber no pvlib
\item Modelo de sujidade HSU
\item RdTools SRR; Deceglie et al.\ (2018), \emph{IEEE JPV} 8(2)
\item IEA PVPS T13-21:2022, perdas por sujidade
\item Limpeza por chuva, \emph{Solar RRL} (2024)
\item Hansen et al.\ (2023), reconciliação espacial de previsões eólicas, \emph{Wind Energy}
\item Reconciliação orientada a valor: \url{https://arxiv.org/abs/2501.16086}
\end{enumerate}}
